
\chapter{Introduction}\label{introduction}

This research presents new software to increase user input efficiency
for non-QWERTY lexical input sets (\emph{\textbf{lexica}}), such as
accented Latin letters, non-Latin letters, and Chinese words and
phrases, all of which are represented in the Unicode (Allen et al.,
2012a) character set.

This software developed in this research, \mulex,
facilitates non-QWERTY input. \mulex is designed to
optimize the number of Unicode characters sent to computer apps (e.g.,
word processors) for the number of user gestures on keyboards, mice, and
specialized input peripherals.

A unique feature of \mulex is its flexibility. It can be
used for input sets as simple as ``Spanish Letters'' and as complex as
``Chinese Phrases''.

Chapter 2 discusses the types of lexica, the peripherals used for
lexical input, defines gesture as it applies to peripherals, and posits
gestural efficiency as a measure for input systems.

Chapter 3 presents an in-depth case-study and description of a
\mulex application for the input of Chinese characters,
words, and phrases.

Chapter 3 uses the Chinese input case study and application description
to discuss \mulex design concepts including associating
semantics with alternative actuations, associating visual feedback with
actuations, and arranging sequenced choice sets for optimizing gestural
efficiency.

Chapter 4 generalizes application design in terms of the design
concepts.

Chapter 5 discusses an app typology for lexical input apps and
associating it with characteristics of lexica.

